\subsection{复局部凸空间}

说复向量空间 E 的子集 A 是均衡的，意思是：对所有 \(x \in A\)，有 \(\rho x \in A\)（对 \(0 \leqslant \rho \leqslant 1\)），且 \(e^{i\vartheta}x \in A\) 对所有实数 \(\vartheta\) 成立。

我们说 E 的一个集合 A 是凸的，如果它在 E 底下的实空间 \(E_{0}\) 中是凸的。为使 E 的凸集 \(A \neq \varnothing\) 是均衡的，只需 \(e^{i\vartheta}A \subset A\) 对所有实数 \(\vartheta\) 成立；因为这首先蕴含 -A = A；由 A 是凸的，我们看出 0 属于 A，从而 \(\rho A \subset A\)（对 \(0 \leqslant \rho \leqslant 1\)）。

设 E 是一个复拓扑向量空间。包含 E 的集合 A 的最小均衡凸集（相应地，闭均衡凸集）称为 A 的均衡凸包（相应地，闭均衡凸包）；A 的闭均衡凸包是 A 的均衡凸包的闭包。后者是诸集 \(e^{i\theta}A\) 之并的凸包；因此我们可以把它定义为诸线性和 \(\sum_{i}\lambda_{i}x_{i}\) 所成的集，其中 \((x_{i})\) 是 A 的点的任意有限族，而 \((\lambda_{i})\) 是满足 \(\sum_{i}|\lambda_{i}| \leqslant 1\) 的复数族。若 A 是预紧的，则其均衡包也是预紧的（I, p.~6, 命题 3）。

我们称复拓扑向量空间 E 是局部凸的，如果 E 底下的实拓扑向量空间 \(E_{0}\) 是局部凸的，也就是说，如果 E 中 0 的每个邻域都包含 0 的一个凸邻域；E 上的一个拓扑 T 称为局部凸的，如果它与 E 的（关于 C 的）向量空间结构相容，并且赋有拓扑 T 的 E 是局部凸的。由于在这种情形，0 的每个闭凸邻域 V 都包含一个均衡邻域 W（I, p.~7, 命题 4），我们看出 V 还包含 U，即 W 的闭均衡凸包；换言之，0 的均衡闭凸邻域构成 E 中 0 的基本邻域系，它在每个比率 \(\neq 0\) 的位似下不变。

反之，设 E 是一个复向量空间，并设 S 是 E 上由吸收的均衡凸集组成的一个滤子基。这时我们知道（II, p.~23, 命题 1），由 S 中的集经比率为 \textgreater{} 0 的位似变换得到的诸集所成之集 B，构成 E 底下的实向量空间 \(E_{0}\) 上某个局部凸拓扑 T 的 0 的基本邻域系。此外，由于 B 中的集都是均衡的，它们在每个位似 \(x \mapsto e^{i\vartheta}x\) 下不变，这表明 T 与 E 的（C 上的）向量空间结构相容（I, p.~7, 命题 4）。

复向量空间 E 上的每个局部凸拓扑都可以由一组半范数定义，因为 0 的一个开均衡凸邻域的度规是 E 上的一个半范数。

对实局部凸空间详细给出的、载于 II, p.~25 至 36 的概念与结果，都可毫无修改地推广到复局部凸空间，只需把对称凸集换成均衡凸集。

复局部凸空间若可度量化且完备，就称为 Fréchet 空间。

